\section{Einleitung}

\gls{ai} hat sich zu einem zentralen Forschungs- und Anwendungsgebiet der Informatik entwickelt, und die jüngsten Fortschritte in der Sprachverarbeitung beruhen vor allem auf der \gls{transformer}-Architektur und großen Sprachmodellen \cite{vaswani2017attention, radford2018improving}. \gls{llm} erreichen bei vielen sprachbezogenen Aufgaben hohe Leistung, ihr internes Verhalten bleibt jedoch häufig undurchsichtig. Die \gls{mechanisticinterpretability} setzt genau hier an: Statt nur zu beobachten, was ein Modell ausgibt, untersucht sie, wie und warum eine Ausgabe im Inneren des Modells entsteht \cite{nanda2023comprehensive, elhage2021mathematical}.

Als Fallstudie dient in dieser Arbeit die \gls{sentimentanalysis}. Sentiment lässt sich gut mechanistisch untersuchen, weil es sich an einzelnen sentimenttragenden Wörtern festmachen lässt und sich positive und negative Beispiele mit nahezu identischem Kontext gegenüberstellen lassen. Untersucht wird das offene Modell Pythia-410M von EleutherAI \cite{biderman2023pythia}, ausschließlich mit PyTorch und HuggingFace, ohne Spezialbibliotheken.

\subsection{Hypothesen}

Jede Hypothese ist mit einem konkreten Kriterium versehen, das angibt, wann sie als bestätigt gilt.

\begin{description}[style=nextline, leftmargin=0pt, font=\bfseries, itemsep=0.4em]
  \item[H1 (klare Stimmungswörter)] Eindeutig polare Wörter wie „wonderful“ oder „delicious“ werden früh und direkt am Token verarbeitet. Bestätigt, wenn die Polarität bereits in frühen bis mittleren Schichten ablesbar ist \emph{und} ein Eingriff am Stimmungswort die Vorhersage wiederherstellt.
  \item[H2 (Lokalisierbarkeit)] Die kausal verantwortliche Verarbeitung lässt sich auf eine bestimmte Schicht und Token-Position eingrenzen, statt über das gesamte Netz verteilt zu sein. Bestätigt, wenn eine scharf begrenzte Stelle das korrekte Verhalten nahezu vollständig wiederherstellt.
  \item[H3 (Ironie und Sarkasmus)] Das Modell erkennt keine Ironie, sondern verarbeitet nur die wörtliche Oberflächen-Polarität. Bestätigt, wenn das Modell durchgehend der wörtlichen Lesart folgt und die gemeinte ironische Polarität an keiner Schicht linear ablesbar ist.
\end{description}

